\section{Introduction}

Maintainers, vulnerability scanners, and downstream users often need a concrete answer to a simple question: which released versions are affected by this CVE? This answer drives upgrade guidance, dependency triage, emergency patching, and risk assessment. A scanner that misses an affected release leaves users exposed, while a scanner that marks safe releases as affected creates unnecessary upgrade work. Public CVE metadata and advisory ranges are useful starting points, but they are not always aligned with the actual release history of a project. The affected-version identification task therefore asks for a repository-grounded output: for each CVE, predict the affected release tags or affected versions in the official release universe~\cite{howfar2025}. The desired output is a complete released-version set, not a best-effort advisory interval.

This task is hard because vulnerability fixes rarely map cleanly to a linear version interval. A fixing commit may contain the root-cause fix together with tests, wrappers, refactoring, or unrelated cleanup. The same fix may be backported to several release branches with different commit hashes, or applied as an equivalent fix that does not match the original patch text. Some CVEs require multiple fixes, receive partial fixes, or cross branch boundaries through cherry-picks and merges. Add-only fixes are also common: the patch adds a guard or validation check, but the vulnerable old statement is not deleted. Refactoring can move the same vulnerability condition across files, and a later release may contain similar code but a different guard or reachable sink. In these cases, the old code, the fixing commit, and the released-version set are connected by history, branch structure, and vulnerability conditions, not by a single textual range.

Existing work gives important evidence but does not fully close this gap. The affected-version benchmark by Chen et al. separates vulnerability-level exact match from version-level precision and recall, showing that partial success on individual tags does not imply recovery of the complete released-version set~\cite{howfar2025}. Affected-version methods such as V-SZZ, developer-log based analysis, and Vercation connect fixing commits, vulnerable code, and release versions more directly than advisory-only reasoning~\cite{vszz2022,tdsc2024,vercation2026}. SZZ, LLM-assisted SZZ, agent-assisted SZZ, MAS-SZZ, and temporal-graph search improve how vulnerable statements and vulnerability-inducing commits are found~\cite{llm4szz2025,agentszz2026,masszz2026,beyondblame2026}. Matching-based and recurring-vulnerability detectors provide scalable evidence about vulnerable code, patch signatures, and reused vulnerable components~\cite{fire2024,movery2022,vuddy2017,vulture2025}. However, a bug- or vulnerability-inducing commit is a historical event, while an affected-version set is a release-tag-level vulnerability state. Code presence and patch resemblance are also evidence, not final affected-version decisions.

\toolname{} moves from single-BIC localization to vulnerability-state reconstruction for affected-version identification. Instead of forcing a CVE history into one inducing commit and one fixing commit, \toolname{} represents the history as role-labeled history events. These roles include introduction, prerequisite, refactoring, fix-series, boundary, unrelated, and uncertain events. The roles define how vulnerability state appears, persists, moves, and disappears over the Git DAG. The propagated state is then projected to official release tags to produce the affected-version set.

\toolname{} takes CVE metadata, fixing commits, and repository history as input. It first organizes attacker-relevant vulnerability conditions in an attacker-condition knowledge graph, including attacker-controlled input, trigger condition, exploit precondition, reachable sink, missing guard, unsafe state, and impact path. These conditions guide root-cause reasoning, pre-fix anchor selection, and history-event judgment. \toolname{} then reconstructs evidence-constrained history events, judges boundary events, propagates vulnerability state over the Git DAG, and performs release tag projection. The method does not ask an LLM to issue independent verdicts for every release tag. Semantic judgment is scoped to repository evidence, candidate history events, and boundary decisions. Finally, vulnerability-pattern-driven self-evolution accumulates reviewed evidence, failure modes, anchor-selection outcomes, Judge decisions, and reusable priors across related CVEs. This feedback is evidence-gated: it is tied to provenance and verified evidence, not free-form model memory.

This paper makes the following contributions.
\begin{itemize}
  \item \textbf{From BIC to vulnerability state for affected-version identification.}
  We formulate affected-version identification as vulnerability-state reconstruction rather than single-BIC localization. \toolname{} models a CVE history as role-labeled history events, including introduction, prerequisite, refactoring, fix-series, boundary, unrelated, and uncertain events. These event roles define how vulnerability state propagates over the Git DAG before being projected to release tags.

  \item \textbf{Attacker-condition knowledge graph.}
  We introduce an attacker-condition knowledge graph that structures attacker-controlled inputs, trigger conditions, exploit preconditions, reachable sinks, missing guards, unsafe states, and impact paths. These conditions guide root-cause reasoning, pre-fix anchor selection, and history-event judgment with source-level evidence.

  \item \textbf{Vulnerability-pattern-driven self-evolution.}
  We design an evidence-gated self-evolution mechanism that accumulates reviewed evidence, failure modes, anchor-selection outcomes, Judge decisions, and reusable priors across CVEs sharing repositories, CWEs, root-cause patterns, patch patterns, or attack behaviors. The accumulated patterns are reused as provenance-backed priors for later analyses.
\end{itemize}

We evaluate \toolname{} on the public affected-version benchmark used by prior work and report vulnerability-level Exact Accuracy, No-Miss Rate (NMR), and version-level micro precision, recall, and F1~\cite{howfar2025}. \need{Fill in Exact Accuracy, NMR, and version-level micro F1 after the full benchmark run.} The evaluation also studies where failures arise across root-cause modeling, anchor selection, history reconstruction, boundary judgment, Git-DAG propagation, and release tag projection. We further analyze sensitivity to patch type, modification scope, and branch context, and evaluate the attacker-condition knowledge graph and vulnerability-pattern-driven self-evolution through ablations. \need{Report the attacker-condition knowledge graph ablation after the full benchmark run.} \need{Report the vulnerability-pattern-driven self-evolution ablation after the full benchmark run.} The rest of this paper is organized as follows. Section~\ref{sec:background} defines the task and problem boundaries. Section~\ref{sec:approach} presents the \toolname{} approach. Section~\ref{sec:evaluation} describes the evaluation design. Section~\ref{sec:threats} discusses limitations and threats to validity. Section~\ref{sec:related} reviews related work.
